\documentclass[11pt]{article}
\usepackage[T1]{fontenc}
\usepackage{lmodern}
\usepackage[margin=.92in]{geometry}
\usepackage{amsmath,amssymb,amsthm}
\usepackage{booktabs,listings,microtype,needspace}
\usepackage[hidelinks]{hyperref}
\hypersetup{pdftitle={Table-Free Evaluation of Cayley--Dickson Basis Signs in Constant Word Time},
  pdfauthor={Alika Perdue}}
\urlstyle{same}

\newtheorem{theorem}{Theorem}
\newcommand{\ctz}{\operatorname{ctz}}
\newcommand{\lowbit}{\operatorname{lowbit}}
\newcommand{\pc}{\operatorname{popcount}}
\newcommand{\bor}{\mathbin{\vert}}
\newcommand{\band}{\mathbin{\&}}
\lstset{basicstyle=\small\ttfamily,columns=fullflexible,keepspaces=true,
  showstringspaces=false,xleftmargin=1em,aboveskip=.65\baselineskip,
  belowskip=.65\baselineskip}

\title{Table-Free Evaluation of Cayley--Dickson\\Basis Signs in Constant Word Time}
\author{Alika Perdue\\[.25em]\small Yotta Space Corp, Vancouver, BC, Canada}
\date{CayleyKernel technical report}

\begin{document}
\maketitle

\begin{abstract}
In a real Cayley--Dickson algebra, the product of two basis elements has an
index given by bitwise XOR and a sign determined by the doubling convention.
This report derives the low-bit mask evaluator used by CayleyKernel from the
ordinary doubling product. For distinct nonzero indices, the lowest set bits of the
operands and of their XOR identify a position below which the doubling
recurrence reduces to an identity or a diagonal product. The higher positions
contribute the parity of the operand OR. The resulting unsigned evaluator
uses constant word operations and constant auxiliary words per sign when both
indices fit in one word and population count is a unit-cost primitive.
The report relates the evaluator to Ren and Zhao's twist formula and
documents the Python, Rust, CUDA, and HIP implementations at a pinned
CayleyKernel revision. An executable C listing and independent recurrence
checks make the sign evaluation reproducible. The word bound applies to an
individual sign. Arbitrary-precision indices require width-dependent work,
and direct dense multiplication still evaluates a quadratic number of terms.
\end{abstract}

\section{Introduction and related work}
A direct Cayley--Dickson product repeatedly asks whether to add or subtract a
coefficient product. A sign table answers that question by lookup. At
dimension \(N\), a full table has \(N^2\) entries, while a direct descent
through the doubling construction can recover one sign from the two indices
alone. The distinction between a sign query and the full coefficient
contraction determines which costs a table-free sign rule can remove.

Albuquerque and Majid formulated the octonions and higher Cayley--Dickson
algebras using a twist of a group algebra \cite{albuquerque}. Ren and Zhao
gave an explicit general twist formula \cite{ren,renv2}. Maher Ben Abdessalem
developed a table-free constant-word evaluator by collapsing a descent to a
lowest-bit test and population-count parity. CayleyKernel credits his work
for its evaluator \cite{abdessalem,cayley}. We cite the author-posted
manuscript dated 24 September 2026, whose Section~6 states the evaluator.
This technical report documents CayleyKernel's implementation of the known
sign law. It presents a derivation of the implemented mask form, relates it
to the published formula, and provides reproducible checks for a pinned
software revision.

\section{Doubling and basis signs}
Set \(A_0=\mathbb R\), with conjugation the identity, and let
\(A_n=A_{n-1}\oplus A_{n-1}\). We use ordinary real doubling
\begin{equation}
 (a,b)(c,d)=(ac-\bar d b,\ da+b\bar c),
 \qquad \overline{(a,b)}=(\bar a,-b).
 \label{eq:doubling}
\end{equation}
At level \(n\geq 1\), write \(m=2^{n-1}\) and embed the lower-level
basis by \(e_i\mapsto(e_i,0)\) and \(e_{m+i}=(0,e_i)\), where
\(0\leq i<m\). The label \(e_0\) denotes the unit at every level.
Conjugation gives \(\bar e_b=\eta(b)e_b\), where \(\eta(0)=1\)
and \(\eta(b)=-1\) for \(b>0\). This follows inductively because the
lower half inherits conjugation and the upper half is negated.

Write \(r=\alpha m+a\) and \(c=\beta m+b\), with
\(\alpha,\beta\in\{0,1\}\). Substituting the basis embeddings into
\eqref{eq:doubling} leaves one product in each case. The table uses
\(s_{n-1}\) for the lower-level sign.
\begin{center}
\begin{tabular}{@{}c l l@{}}
\toprule
\((\alpha,\beta)\) & Surviving lower-level product & Sign \(s_n(r,c)\)\\
\midrule
\((0,0)\) & \(e_a e_b\) in the lower half & \(s_{n-1}(a,b)\)\\
\((0,1)\) & \(e_b e_a\) in the upper half & \(s_{n-1}(b,a)\)\\
\((1,0)\) & \(e_a\bar e_b\) in the upper half & \(\eta(b)s_{n-1}(a,b)\)\\
\((1,1)\) & \(-\bar e_b e_a\) in the lower half & \(-\eta(b)s_{n-1}(b,a)\)\\
\bottomrule
\end{tabular}
\end{center}
At level zero \(e_0e_0=e_0\). By induction the surviving product has
lower index \(a\oplus b\) in either order and the output half is
\(\alpha\oplus\beta\), where \(\oplus\) is bitwise XOR. Thus
\begin{equation}
 e_r e_c=s_n(r,c)e_{r\oplus c},\qquad s_n(r,c)\in\{+1,-1\}.
 \label{eq:twist}
\end{equation}
The XOR selects the output basis element. The sign is a scalar twist over
\(\mathbb R\). Arithmetic modulo two is used for bits and sign
exponents. Associativity would require
\(s(r,c)s(r\oplus c,t)=s(c,t)s(r,c\oplus t)\) for all triples.
The octonions are nonassociative, so the twist in this family need not
satisfy the associative 2-cocycle identity \cite{albuquerque}.

The recurrence gives the following identities by induction on \(n\).
When \(r=0\), the \((0,0)\) and \((0,1)\) cases reduce to a
lower-level product with a zero index. When \(c=0\), the \((0,0)\)
and \((1,0)\) cases do the same, with \(\eta(0)=1\). Thus both
identity cases give \(+1\), including \((0,0)\). A nonzero diagonal
in the lower half inherits its sign. In the upper half its sign is
\(-\eta(a)s_{n-1}(a,a)\), which is \(-1\) for \(a=0\) and
\(-1\) for \(a>0\) by the induction hypothesis.

For distinct nonzero indices in the lower half, antisymmetry is inherited.
For a mixed pair with the lower-half index first, that index is \(a>0\).
Its sign is \(s_{n-1}(b,a)\) and the reversed sign is
\(\eta(a)s_{n-1}(b,a)=-s_{n-1}(b,a)\). For an upper-half pair the
suffixes are distinct. If both are nonzero, the factors
\(-\eta(a)=-\eta(b)=1\) leave antisymmetry to the lower-level signs.
If one suffix is zero, both lower products have sign \(+1\) and the
two conjugation factors have opposite signs. This proves
\(s_n(r,c)=-s_n(c,r)\) for distinct nonzero inputs. One extra leading
zero bit in both inputs uses the \((0,0)\) case and preserves the sign,
so we omit the level subscript when the indices fit.

The table is also an independent sign algorithm. Start with a sign
accumulator of \(+1\) and process bits from highest to lowest. At each
position strip both top bits, negate the accumulator if the row bit was
set and the column suffix is nonzero, and negate again if both top bits
were set. If the column bit was set, swap the two remaining suffixes.
The final accumulator is the sign because \(s_0(0,0)=1\). The oracle
in Appendix~\ref{app:c} implements this descent.

\section{From the recurrence to one word}
For a nonzero integer \(x\), \(\ctz(x)\) is the position of its least
significant set bit, numbered from zero. We write \(r_i,c_i\) for the
bits at position \(i\) and \(u=r\bor c\).

In a descent from the highest bit, as long as the two lower suffixes
are nonzero and distinct, antisymmetry lets us account for a swap
by flipping the sign. The four doubling cases then reduce to a
single OR bit at each position. The position \(k\) below marks where
one suffix first becomes zero or the two become equal. At that point,
the identity or diagonal rule evaluates the lower product directly.

\Needspace{13\baselineskip}
\begin{theorem}[Mask parity for ordinary doubling]
\label{thm:mask}
Let \(0\leq r,c<2^n\). A zero operand gives sign \(+1\), and
an equal nonzero pair gives sign \(-1\). Otherwise define
\[
 h=\ctz(r),\quad v=\ctz(c),\quad d_0=\ctz(r\oplus c),
 \qquad k=\max(h,v,d_0).
\]
Set \(q=1-c_k\) if \(h>v\), \(q=r_k\) if \(v>h\), and
\(q=c_k\) if \(h=v\). Then
\begin{equation}
 s(r,c)=(-1)^{\,q\oplus U},\qquad
 U=\bigoplus_{i>k}(r_i\lor c_i).
 \label{eq:upper}
\end{equation}
Equivalently, the mask evaluator in Listing~\ref{lst:mask} returns
\(s(r,c)\). It uses \(O(1)\) word operations and \(O(1)\) auxiliary
words per sign when the two indices fit in one unsigned word and word
population count is a unit-cost primitive.
\end{theorem}

\begin{proof}
The identity and diagonal cases were proved from doubling and are handled
before any trailing-zero counts or shifts. For the remaining inputs, if
\(h\ne v\), exactly one operand has a set bit at \(\min(h,v)\) and
both have zero bits below it. Hence \(d_0=\min(h,v)\). If \(h=v\),
the operands agree up to and including bit \(h\), so \(d_0>h=v\).
In each case the maximum \(k\) is unique among the three positions.

At a position \(i>k\), each suffix below \(i\) contains its first set
bit, and the two suffixes differ at \(d_0<i\). Write them in original
order as \((a,b)\). They are nonzero and distinct. For \((0,0)\) the
sign is \(s(a,b)\). For \((0,1)\) the swapped product contributes
\(s(b,a)=-s(a,b)\). For \((1,0)\) the factor \(\eta(b)=-1\)
gives \(-s(a,b)\). For \((1,1)\) the factor \(-\eta(b)=1\)
leaves \(s(b,a)=-s(a,b)\). The normalized lower pair therefore
stays in its original order and each of the three nonzero bit pairs
contributes one flip. The exponent of these flips is \(U\).

At \(i=k\), write \(a=r\bmod 2^k\) and \(b=c\bmod 2^k\).
The first-set and first-differing positions give the following cases.
\begin{center}
\begin{tabular}{@{}llll@{}}
\toprule
Order & Suffixes & Bits at \(k\) & Local exponent \(q\)\\
\midrule
\(h>v=d_0\) & \(a=0,\ b>0\) & \(r_k=1\) & \(1-c_k\)\\
\(v>h=d_0\) & \(a>0,\ b=0\) & \(c_k=1\) & \(r_k\)\\
\(d_0>h=v\) & \(a=b>0\) & \(r_k\ne c_k\) & \(c_k\)\\
\bottomrule
\end{tabular}
\end{center}
If \(h>v\), \(c_k=0\) gives \(\eta(b)s(0,b)=-1\), and
\(c_k=1\) gives \(-\eta(b)s(b,0)=+1\).
If \(v>h\), \(r_k=0\) gives \(s(0,a)=+1\), and
\(r_k=1\) gives \(-\eta(0)s(0,a)=-1\).
If \(h=v\), the bits \((1,0)\) give \(\eta(a)s(a,a)=+1\)
and \((0,1)\) gives \(s(a,a)=-1\).
This proves \eqref{eq:upper}. Exchanging the operands preserves \(k\)
and \(U\). It interchanges the first two cases, replacing \(1-c_k\)
by \(c_k\) or \(r_k\) by \(1-r_k\). In the third case it replaces
\(c_k\) by \(r_k=1-c_k\). The sign therefore reverses.

For the mask form, work in unsigned words modulo \(2^w\). In
\(-x=\mathord{\sim}x+1\), the trailing zeros of \(x\) reappear as
zeros and its first set bit reappears as one, while all higher bits are
complemented. Consequently \(x\band(-x)=2^{\ctz(x)}\) for
\(x\ne0\). For a one-hot mask \(M=2^k\), \(-M=2^w-2^k\)
has ones from \(k\) through \(w-1\). Hence \(u\band(-M)\)
retains the suffix at and above \(k\). In every case \(u_k=1\),
so \(P=\pc(u\gg k)\bmod 2=1\oplus U\).
Here \(\gg\) is a logical right shift.

Let \(R=\lowbit(r)\), \(C=\lowbit(c)\), and
\(d=\lowbit(r\oplus c)\). Comparing these masks compares their
bit positions. The test \(C<R\) selects \(h>v\) and
\(\mathord{\sim}c\band R=qR\). If that test fails, \(d<C\)
selects \(v>h\) and \(r\band C=qC\). The remaining case
\(h=v\) uses \(c\band d=qd\). Thus the final mask is
\((u\band(-M))\oplus(qM)\), where \(M=R,C,d\), respectively.
Population-count parity is linear under XOR, so the mask parity is the
\emph{positive-sign flag} \(p=P\oplus q\). Equation~\eqref{eq:upper}
gives \(s=2p-1=(-1)^{1\oplus p}\). The listing uses a bounded number
of word operations with one population count and a bounded number of
temporary words, proving the word-cost claim.
\end{proof}

\begin{lstlisting}[float=tbp,caption={Unsigned word evaluator. Negation and NOT act on the full word.},label={lst:mask}]
basis_sign(r, c):
    if r == 0 or c == 0:
        return +1
    if r == c:
        return -1

    d = lowbit(r XOR c)
    R = lowbit(r)
    C = lowbit(c)
    u = r OR c
    if C < R:
        mask = ((NOT c) AND R) XOR (u AND (-R))
    else if d < C:
        mask = (r AND C) XOR (u AND (-C))
    else:
        mask = (c AND d) XOR (u AND (-d))
    p = popcount(mask) AND 1
    return 2*p - 1
\end{lstlisting}

\subsection{Shifted population count and the general twist formula}
The shifted-popcount form uses \(P\) directly. In the third case
\(r_k\ne c_k\), so testing whether bit \(k\) of \(r\) is clear
also obtains \(q=c_k\). These are precisely the three tests in the
Rust evaluator, ordered as \(h>v\), \(v<d_0\), and the remaining
case. Under the trailing-zero identities, \(v<d_0\) is equivalent to
\(h=v\). With fixed limbs, the first retained limb is shifted locally
and the popcount parities of all retained limbs are XORed. The result is
the same \(P\) because the limbs partition the retained bits.

Ren and Zhao's arXiv v2, Theorem~1 and equation (4.1), use
\(e_Ae_B=(-1)^{\sigma(A,B)}e_{A\oplus B}\) \cite{renv2}.
Their equation (2.1) is \eqref{eq:doubling} in \(a+bg_n\)
notation. Multiplication on the right by the newest generator
\((0,1)\) sends \((a,0)\) to \((0,a)\). Their left-associated,
ascending generator labels thus agree inductively with half-split indexing.
Their degree is \(\ctz\), their \(l\) is our \(k\), and their
sum of bitwise OR terms from \(l\) upward equals \(P\) modulo two.
Their three additions to the sum are \(c_k\), \(1\oplus r_k\),
and \(r_k\) for \(h>v\), \(h<v\), and \(h=v\), respectively.
Each equals \(1\oplus q\), using \(r_k\ne c_k\) in the last
case. Thus \(\sigma=P\oplus(1\oplus q)=1\oplus p\). Equation
(4.1) assigns exponent zero on the identity and one on the nonzero
diagonal, also agreeing with the evaluator.

\subsection{Worked products}
The following calculations show three-bit words with the least significant
bit on the right. Unsigned negatives and complements are truncated to three
bits for display. For each pair, the mask is formed from the retained
operand OR and the selected correction bit.

\begin{description}
\item[\((2,1)=(010,001)_2\).]
Here \((R,C,d)=(010,001,001)_2\) and \(u=011_2\).
The test \(C<R\) selects \(M=R=010_2\), with \(-M=110_2\).
Then \(u\band(-M)=010_2\) and \((\mathord{\sim}c)\band R=010_2\).
Their XOR is \(000_2\), so \(p=0\). The index is
\(010_2\oplus001_2=011_2\), giving \(-e_3\).
\item[\((1,6)=(001,110)_2\).]
Here \((R,C,d)=(001,010,001)_2\) and \(u=111_2\).
The test \(d<C\) selects \(M=C=010_2\), with \(-M=110_2\).
Then \(u\band(-M)=110_2\) and \(r\band C=000_2\).
Their XOR is \(110_2\), which has even parity \(p=0\).
The index is \(001_2\oplus110_2=111_2\), giving \(-e_7\).
\item[\((5,7)=(101,111)_2\).]
Here \((R,C,d)=(001,001,010)_2\) and \(u=111_2\).
The last branch selects \(M=d=010_2\), with \(-M=110_2\).
Then \(u\band(-M)=110_2\) and \(c\band d=010_2\).
Their XOR is \(100_2\), so \(p=1\).
The index is \(101_2\oplus111_2=010_2\), giving \(+e_2\).
\item[Reversal.]
Reversing the first pair gives \((1,2)=(001,010)_2\),
with \((R,C,d)=(001,010,001)_2\) and \(u=011_2\).
Here \(M=C=010_2\), \(-M=110_2\), and the mask is
\(010_2\oplus000_2=010_2\), so \(p=1\) and the product is
\(+e_3\).
\item[Identity and diagonal.]
For \((0,5)\), the guard returns \(+1\) without a
mask and the XOR gives index 5. For \((5,5)\), the guard returns
\(-1\) without a mask and the XOR gives index 0. The pair \((0,0)\)
uses the identity guard and gives \(+e_0\).
\end{description}

\section{Cost model and dense multiplication}
The table compares sign evaluators in the single-word model of
Theorem~\ref{thm:mask}. The table row assumes directly addressable
entries and excludes the cost of building the table. The doubling row
uses an iterative scan over the \(n\) bit positions.

\begin{center}
\begin{tabular}{@{}lll@{}}
\toprule
Evaluator & Worst-case word work per sign & Sign storage or workspace\\
\midrule
Full precomputed table & \(\Theta(1)\) & \(\Theta(N^2)\) stored bits\\
Iterative doubling & \(\Theta(n)=\Theta(\log_2 N)\) & \(O(1)\) auxiliary words\\
Low-bit mask & \(\Theta(1)\) & \(O(1)\) auxiliary words\\
\bottomrule
\end{tabular}
\end{center}

For unbounded \(n\)-bit indices a direct implementation of the mask
rule uses \(O(n)\) bit work or \(O(\lceil n/w\rceil)\) machine-word
work with word population count. The parity alone can scan every limb
of the retained suffix. For a bit-probe lower bound, fix \(r=1\),
\(c_0=0\), and \(c_1=1\), leaving \(c_i\) free for \(2\leq i<n\).
Then \(h=d_0=0\), \(v=k=1\), \(q=r_1=0\), and
\(s(1,c)=(-1)^{\bigoplus_{i=2}^{n-1}c_i}\). Any exact bit-probe
algorithm must inspect every free bit in the worst case because changing
an unprobed bit reverses the answer. The resulting \(\Omega(n)\)
bound is for bit probes. Python integer and Rust \texttt{BigUint}
operations have width-dependent costs.

For \(x=\sum_r x_r e_r\) and \(y=\sum_c y_c e_c\), bilinearity
gives \(xy=\sum_{r,c} x_r y_c s(r,c)e_{r\oplus c}\).
Writing \(xy=\sum_k z_k e_k\), each row \(r\) contributes to
output \(k\) at the unique column \(c=k\oplus r\). Therefore
\begin{equation}
 z_k=\sum_{r=0}^{N-1}s(r,k\oplus r)x_r y_{k\oplus r}.
 \label{eq:dense}
\end{equation}
The direct dense kernel uses \(\Theta(N)\) terms per output and
\(\Theta(BN^2)\) total arithmetic work for \(B\) independent products.
Device threads parallelize outputs and batch items. The direct scheme
stores \(\Theta(BN)\) input and output coefficients and needs no
\(N\)-by-\(N\) sign table. This arithmetic work bound describes the
direct kernel. Elapsed time also depends on scheduling, arithmetic
throughput, memory traffic, and any interface and transfer costs included
in the measurement.

\section{CayleyKernel implementation}
All source descriptions refer to commit \texttt{926f8e8} of
CayleyKernel \cite{cayley}. The software is distributed under
Apache License 2.0.

\paragraph{Python.}
\path{CayleyKernel/_cpu.py} implements the low-bit cases in
\texttt{\_python\_twist} using Python integers and \texttt{bit\_count()}.
Its locals \texttt{v} and \texttt{h} correspond to \(R\) and \(C\),
respectively. Python's unbounded complement and negation yield the same
selected bits because the result is intersected with finite nonnegative
masks. The public \texttt{twist(r,c)} takes nonnegative zero-based
indices and returns \(\pm1\). The public \texttt{basis(r,c)} takes
nonzero signed one-based labels. For example, \texttt{twist(2,1)}
returns \(-1\) and \texttt{basis(3,2)} returns \(-4\), the label
for \(-e_3\).

\paragraph{Rust.}
In \path{rust/src/lib.rs}, \texttt{twist} dispatches on the
larger operand bit length. Up to 64 bits it uses \texttt{twist\_u64}.
From 65 through 128 bits it uses two stack-allocated \texttt{u64}
limbs per operand, and from 129 through 256 bits it uses four.
The fixed-limb path finds the first nonzero limb of each operand and
their XOR, selects the correction bit, and accumulates retained-limb
parity. Above 256 bits \texttt{twist\_big} uses \texttt{BigUint}
trailing zeros, XOR, OR, shift, and population count. Both the Python
and Rust interfaces support nonnegative indices beyond machine-word
width, subject to available resources.

\paragraph{CUDA and HIP.}
The \texttt{cayley\_sign} bodies in \path{cuda/cuda_kernel.cu} and
\path{cuda/amd_kernel.cu} use the mask form and \texttt{\_\_popc}.
The guard in each ternary selects a popcount input of one for a zero
operand and zero for a nonzero diagonal. The bodies return flag 1 to
add a coefficient term and flag 0 to subtract it. They compute low-bit
masks directly and contain no \texttt{\_\_ffs} call. The bodies have
a bounded operation count on their fixed-width operands. Their source
uses conditional expressions, and generated branch behavior, instruction
latency, and hardware throughput require separate evidence.

The device evaluators take nonnegative signed 32-bit \texttt{int}
indices. For such inputs \texttt{r}, \texttt{c}, their XOR, and
the extracted low bits are at most \(2^{31}-1\), so negating them
fits the signed type. Dense paths use \(0\leq r,c<\texttt{dim}\)
with \texttt{dim} also passed as a signed \texttt{int}. A positive
power-of-two dimension in that type is at most \(2^{30}\), an ABI
ceiling rather than a tested capacity. Memory and launch constraints
apply separately. The appendix's unsigned C reference, in contrast,
accepts all \(0\leq r,c<2^{32}\) as individual indices.

\section{Validation and limitations}
\label{sec:validation}
The host checks below ran on 8 November 2026 (9 November UTC) on an
AMD Ryzen 9 7900X with Ubuntu 24.04.5 LTS, Linux
\texttt{7.0.0-42-generic}, and x86-64. The C and C++ compiler was GCC
13.3.0. The Python environment was CPython 3.12.3 with NumPy 2.4.6
and pytest 9.0.3. The Rust extension was built with rustc and Cargo
1.99.1 in release mode from the pinned lockfile (PyO3 0.27.2,
\texttt{num-bigint} 0.4.8). Appendix~\ref{app:reproduce} gives the
commands.

\begin{itemize}
\item The repository's \path{cuda/tools/compare_sign_tables.cpp}
  agreed with independent doubling for all \(16{,}777{,}216\) ordered
  pairs in dimension 4096. Its Mersenne Twister seeded with 2026
  supplied 10,000 pairs masked to 31 bits. The pair, a diagonal,
  and both zero-operand positions gave 40,000 further checks.
  This executable compares an unsigned host translation of the device
  expression with doubling.
\item The appendix C listing agreed with its own iterative doubling
  reference for the same exhaustive range and 50,000 further checks.
  A 32-bit xorshift generator seeded with 2026 supplied 10,000 full-width
  pairs. The check includes each pair, its reversal, a diagonal,
  and both zero-operand positions.
\item The pinned CPU test module \path{tests/core/test_cpu_backend.py}
  passed all 62 tests with both backends available. The companion width
  check passed 603 dispatch-width pairs and 72 limb-boundary pairs per
  backend, plus \((0,0)\). It compares each result to
  \path{tests/support/oracles.py}'s doubling recurrence at the minimum
  level and with seven leading zero bits added. The dispatch cases use
  widths 63, 64, 65, 127, 128, 129, 255, 256, and 257, with 32
  seeded draws and their reversals, two identity orientations, and one
  nonzero diagonal per width. The limb cases place the first set bit or
  first differing bit at positions one below, at, and one above 64, 128,
  192, and 256. The first-set cases include both values of the other
  operand's bit and both operand orders.
\item The CPU class in
  \path{tests/papers/desktop_batch/test_2205_07986v2.py}
  passed all seven tests separately under Python and Rust. Its exhaustive
  test compares equation (4.1), the CPU interfaces, and an independent
  recurrence for all \(1{,}048{,}576\) pairs in dimension 1024.
  Seeded higher-width cases and a check of the final split parameter
  supplement that comparison.
\end{itemize}

\paragraph{Version-specific octonion check.}
Ren and Zhao's arXiv v2 Example 1 gives incorrect degrees for
\((5,6)\). Their definition gives \((h,v,d_0)=(0,1,0)\)
and \(k=1\). Its final product \(-e_3\) agrees with doubling.
The shortcut printed in equation (4.11) uses bit 2 and includes
all three OR terms. For \((1,2)\) it gives exponent
\(r_2+0+1+1+1\equiv1\pmod 2\), whereas equation (4.1) gives
\(1+r_1+\phi(r_1,c_1)+\phi(r_2,c_2)\equiv0\pmod 2\).
Doubling likewise gives \(+e_3\). The repository companion
reproduces this discrepancy for the specified arXiv v2 copy, with
24 mismatches among the 42 ordered pairs of distinct nonzero octonion
indices \cite{renv2,cayley}.

\paragraph{Scope of the final split parameter.}
Replacing the minus in the first component of \eqref{eq:doubling}
by a plus only at the last doubling changes the sign of the
\((1,1)\) case, since that is the only case where the
\(\bar d b\) term survives. The split exponent is therefore
\(\sigma\oplus(r_{n-1}c_{n-1})\), as in Ren and Zhao's
Theorem~2 \cite{renv2}. The evaluator and the implementation
scope of this paper use ordinary doubling.

\paragraph{Evidence scope.}
The finite comparisons corroborate the theorem. The saved
\texttt{2205.07986v2} JSON snapshot records a separate run on
26 October 2026 at commit \texttt{399309c}, including device
coverage in that recorded suite. The fresh runs reported here are
host-only. No CUDA or HIP execution, instruction inspection, or
performance benchmark was part of this run. The repository's
benchmark methodology separates native sign throughput, Python
interface timings, and dense-product timings.

\section*{Declarations}
\paragraph{Author contribution and AI use.}
Alika Perdue developed CayleyKernel with AI assistance, explored sign
patterns in computed basis products, and defined the scope and verification
requirements for this report. OpenAI Codex generated most of the manuscript,
including the derivation and proof, drawing on the author's observations,
the implementation, and the cited mathematical work. Codex also prepared the
companion C and Python checks and ran the verification reported here.

\paragraph{Affiliation and competing interests.}
The author owns Yotta Space Corp and developed CayleyKernel through the
company. The CayleyKernel README carries the company's copyright notice.
% Confirm formal funding status before any submission that requires it.

\appendix

\section{Executable C reference}
\label{app:c}
Listing~\ref{lst:c} is the complete companion file
\path{basis_sign.c}. It uses \texttt{uint32\_t}, requires 32-bit
\texttt{unsigned int}, and uses the GCC/Clang intrinsic
\texttt{\_\_builtin\_popcount}. The doubling reference processes
32 positions and never calls the mask evaluator. The executable
returns a nonzero status on the first mismatch.
\lstinputlisting[language=C,caption={Unsigned mask evaluator and independent test harness.},label={lst:c}]{basis_sign.c}

\section{Reproduction}
\label{app:reproduce}
To build the manuscript, enter its directory and run
\texttt{latexmk -pdf paper.tex}. Keep \path{basis_sign.c} next to
\path{paper.tex} because the listing is read directly from the file.

Set \texttt{paper\_dir} to the absolute path of this manuscript
directory. Run the following in a shell on a Linux host with GCC,
Rust and Cargo, and Python with NumPy and pytest. Use an environment
with the dependencies for the pinned project, for example
\verb|python -m pip install -e '.[test]'| in that checkout.
The Cargo step below builds a CPython extension using the lockfile.
It uses Linux shared-library filenames and places the result on the
test import path.
\Needspace{31\baselineskip}
\begin{lstlisting}
git clone https://github.com/alikaperdue/CayleyKernel.git
cd CayleyKernel
git checkout --detach 926f8e85506187086a8ca2ea3577385f4a578b96
git rev-parse HEAD

check_dir=$(mktemp -d /tmp/cayley-sign.XXXXXX)
c++ -std=c++17 -O3 -Wall -Wextra \
    cuda/tools/compare_sign_tables.cpp -o "$check_dir/compare"
"$check_dir/compare" 4096

cc -std=c11 -O3 -Wall -Wextra -Werror -pedantic \
    "$paper_dir/basis_sign.c" -o "$check_dir/basis_sign"
"$check_dir/basis_sign"

PYO3_PYTHON=$(command -v python) CARGO_TARGET_DIR="$check_dir/target" \
    cargo rustc --locked --release --manifest-path rust/Cargo.toml \
    --features python,pyo3/extension-module --crate-type cdylib
cp "$check_dir/target/release/libcayley_kernel_rust.so" \
    "$check_dir/cayley_kernel_rust.so"
export PYTHONPATH="$PWD:$check_dir"
CAYLEY_TEST_REQUIRE_RUST=1 python -m pytest -q \
    tests/core/test_cpu_backend.py

for backend in python rust; do
    export CAYLEY_KERNEL_CPU_BACKEND="$backend"
    python "$paper_dir/check_widths.py"
    python -m unittest \
        tests.papers.desktop_batch.test_2205_07986v2.Sections -v
done
\end{lstlisting}
The reported Cargo build used \texttt{--offline} with cached crates.
The two files supplied with the manuscript are separate from the pinned
software checkout. The repository's
\href{https://github.com/alikaperdue/CayleyKernel/blob/926f8e85506187086a8ca2ea3577385f4a578b96/docs/papers/2205.07986v2/README.md}{literature companion}
describes the comparison and saved snapshot. Its
\href{https://github.com/alikaperdue/CayleyKernel/blob/926f8e85506187086a8ca2ea3577385f4a578b96/docs/benchmark-methodology.md}{benchmark methodology}
documents the separate timing protocols.

\begin{thebibliography}{9}
\bibitem{albuquerque}
H. Albuquerque and S. Majid.
Quasialgebra Structure of the Octonions.
\emph{Journal of Algebra}, 220(1), 188--224, 1999.
\href{https://doi.org/10.1006/jabr.1998.7850}{doi:10.1006/jabr.1998.7850}.
\href{https://arxiv.org/abs/math/9802116}{arXiv:math/9802116}.

\bibitem{ren}
G. Ren and X. Zhao.
The Explicit Twisted Group Algebra Structure of the Cayley--Dickson Algebra.
\emph{Advances in Applied Clifford Algebras}, 33, article 49, 2023.
\href{https://doi.org/10.1007/s00006-023-01296-6}{doi:10.1007/s00006-023-01296-6}.

\bibitem{renv2}
G. Ren and X. Zhao.
The Twisted Group Algebra Structure of the Cayley--Dickson Algebra.
\href{https://arxiv.org/html/2205.07986v2}{arXiv:2205.07986v2},
13 March 2023. Theorem and equation numbers in this manuscript
refer to this version.

\bibitem{abdessalem}
M. Ben Abdessalem.
An \(O(1)\) Bitwise Evaluator for Cayley--Dickson Sign Structure:
Ordinary, Split, Dual, and Tensor Constructions.
\href{https://www.academia.edu/179628334/An_O_1_Bitwise_Evaluator_for_Cayley_Dickson_Sign_Structure_Ordinary_Split_Dual_and_Tensor_Constructions}{Author-posted preprint}, manuscript dated 24 September 2026.
The posting associates this copy with Figshare DOI
\href{https://doi.org/10.6084/m9.figshare.33705022}{\nolinkurl{10.6084/m9.figshare.33705022}}.
The Figshare version and any journal publication remain unverified.

\bibitem{cayley}
CayleyKernel contributors. \emph{CayleyKernel}. Research software.\\
Commit \href{https://github.com/alikaperdue/CayleyKernel/tree/926f8e85506187086a8ca2ea3577385f4a578b96}
{\nolinkurl{926f8e85506187086a8ca2ea3577385f4a578b96}}.
\href{https://github.com/alikaperdue/CayleyKernel/blob/926f8e85506187086a8ca2ea3577385f4a578b96/LICENSE}{Apache License 2.0}.
\end{thebibliography}
\end{document}
